\section{Exact Gaussian Sampler}
\label{sec:exactgpv}

In this section we prove Lemma~\ref{lem:gpv}. 
As in~\cite{DBLP:conf/stoc/GentryPV08}, the proof consists of two parts. 
In the first we consider the one-dimensional case, and in the second we use
it recursively to sample from arbitrary lattices. Our one-dimensional
sampler is based on rejection sampling, just like the one in~\cite{DBLP:conf/stoc/GentryPV08}.
Unlike~\cite{DBLP:conf/stoc/GentryPV08}, we use the continuous normal distribution as the source distribution
which allows us to avoid truncation, and as a result obtain an exact sample. 
Our second part uses the same recursive routine as in~\cite{DBLP:conf/stoc/GentryPV08},
but adds a rejection sampling step to it in order to take care of the deviation 
of its output from the desired distribution. 

\subsection{The one-dimensional case}

Here we show how to sample from the discrete Gaussian distribution on arbitrary cosets of
one-dimensional lattices. We use a standard rejection sampling procedure (see, e.g.~\cite[Page 117]{DevroyeBook86}
for a very similar procedure).

By scaling, we can restrict without loss of generality to the lattice $\Z$, i.e.,
we consider the task of sampling from $D_{\Z+c,r}$ for a given coset representative $c \in [0,1)$ and parameter $r>0$.
The sampling procedure is as follows. Let $Z_0 = \int_c^\infty \rho_r(x) \mathrm{d} x$, and $Z_1 = \int_{-\infty}^{c-1} \rho_r(x) \mathrm{d} x$.
These two numbers can be computed efficiently by expressing them in terms of the error function. Let $Z=Z_0+Z_1+\rho_r(c)+\rho_r(c-1)$.
The algorithm repeats the following until it outputs an answer: 
\begin{itemize}
\item With probability $\rho_r(c)/Z$ it outputs $c$;
\item With probability $\rho_r(c-1)/Z$ it outputs $c-1$;
\item With probability $Z_0/Z$ it chooses $x$ from the restriction of the continuous normal distribution $D_r$ to the interval $[c,\infty)$. Let $y$
be the smallest element in $\Z+c$ that is larger than $x$. With probability $\rho_r(y)/\rho_r(x)$ output $y$, and otherwise repeat;
\item With probability $Z_1/Z$ it chooses $x$ from the restriction of the continuous normal distribution $D_r$ to the interval $(-\infty,c-1]$. Let $y$
be the largest element in $\Z+c$ that is smaller than $x$. With probability $\rho_r(y)/\rho_r(x)$ output $y$, and otherwise repeat.
\end{itemize}

Consider now one iteration of the procedure. The probability of outputting $c$ is $\rho_r(c)/Z$, that of outputting $c-1$ 
is $\rho_r(c-1)/Z$, that of outputting $c+k$ for some $k \ge 1$ is
\[
\frac{Z_0}{Z} \cdot \frac{1}{Z_0} \int_{c+k-1}^{c+k} \rho_r(x) \cdot \frac{\rho_r(c+k)}{\rho_r(x)}  \mathrm{d} x
= \frac{\rho_r(c+k)}{Z},
\]
and similarly, that of outputting $c-1-k$ for some $k \ge 1$ is 
$\rho_r(c-1-k)/Z$. From this it follows immediately that conditioned on outputting something, the output
distribution has support on $\Z+c$ and probability mass function proportional to $\rho_r$, and is therefore the desired discrete
Gaussian distribution $D_{\Z+c,r}$. Moreover, the probability of outputting something is
\[
\frac{\rho_r(\Z+c)}{Z} = \frac{\rho_r(\Z+c)}{Z_0+Z_1+\rho_r(c)+\rho_r(c-1)} \ge \frac{\rho_r(\Z+c)}{\rho_r(\Z+c)+\rho_r(c)+\rho_r(c-1)}
\ge \frac{1}{2}.
\]
Therefore at each iteration the procedure has probability of at least $1/2$ to terminate. As a result, the probability that the 
number of iterations is greater than $t$ is at most $2^{-t}$, and in particular, the expected number of iterations
is at most $2$.

\subsection{The general case}

For completeness, we start by recalling the $\sampleD$ procedure described in~\cite{DBLP:conf/stoc/GentryPV08}.
This is a recursive procedure that gets as input a basis $\matB=(\vecb_1,\ldots,\vecb_n)$ of an $n$-dimensional lattice $\Lambda=\lat(\matB)$, a parameter $r>0$, 
and a vector $\vecc \in \R^n$, and outputs a vector in $\Lambda+\vecc$ whose distribution is close to that of 
$D_{\Lambda+\vecc,r}$. Let $\widetilde{\vecb_1},\ldots,\widetilde{\vecb_n}$
be the Gram-Schmidt orthogonalization of $\vecb_1,\ldots,\vecb_n$, and let $\overline{\vecb_1},\ldots,\overline{\vecb_n}$ be the
normalized Gram-Schmidt vectors, i.e., $\overline{\vecb_i} = \widetilde{\vecb_i}/\|\widetilde{\vecb_i}\|$. The procedure is the following. 

\begin{enumerate}
\item Let $\vecc_n \leftarrow \vecc$. For $i \leftarrow n,\ldots,1$, do:
\begin{enumerate}
	\item Choose $v_i$ from $D_{\|\widetilde{\vecb_i}\| \Z+\inner{\vecc_i,\overline{\vecb_i}}, r}$ using the exact one-dimensional sampler.
	\item Let $\vecc_{i-1} \leftarrow \vecc_i + (v_i - \inner{\vecc_i, \overline{\vecb_i}}) \cdot \vecb_i / \|\widetilde{\vecb_i}\| - v_i \overline{\vecb_i}$.
\end{enumerate}	
\item Output $\vecv := \sum_{i=1}^n v_i \overline{\vecb_i}$.
\end{enumerate}

It is easy to verify that the procedure always outputs vectors in the coset $\Lambda+\vecc$. Moreover, the probability of
outputting any $\vecv \in \Lambda+\vecc$ is 
\[
\prod_{i=1}^n \frac{\rho_r(v_i)}{\rho_r(\|\widetilde{\vecb_i}\| \Z + \inner{\vecc_i,\overline{\vecb_i}})} = 
\frac{\rho_r(\vecv)}{\prod_{i=1}^n \rho_r(\|\widetilde{\vecb_i}\| \Z + \inner{\vecc_i,\overline{\vecb_i}})},
\]
where $\vecc_i$ are the values computed in the procedure when it outputs $\vecv$.
Notice that by Lemma~\ref{lem:boundonsmoothing} and our assumption on $r$, 
we have that $r \ge \eta_{1/(n+1)}(\|\widetilde{\vecb_i}\| \Z)$ for all $i$.
Therefore, by Lemma~\ref{lem:smoothing}, we have that for all $c \in \R$,
\begin{equation*}%\label{eq:deviationofgaussmass}
\rho_r(\|\widetilde{\vecb_i}\| \Z + c) \in \bracks*{1-\frac{2}{n+2},1} \rho_r(\|\widetilde{\vecb_i}\| \Z).
\end{equation*}

In order to get an exact sample, we combine the above procedure with rejection sampling.
Namely, we apply $\sampleD$ to obtain some vector $\vecv$. We then output $\vecv$
with probability 
\begin{equation}\label{eq:rejsamplingprob}
\frac{\prod_{i=1}^n \rho_r(\|\widetilde{\vecb_i}\| \Z + \inner{\vecc_i,\overline{\vecb_i}})}
{\prod_{i=1}^n \rho_r(\|\widetilde{\vecb_i}\| \Z)} \in \left( \left(1-\frac{2}{n+2}\right)^n,1 \right] \subseteq (e^{-2},1],
\end{equation}
and otherwise repeat. This probability can be efficiently computed, as we will show below.
As a result, in any given iteration the probability of outputting the vector $\vecv \in \Lambda+\vecc$ is 
\[
\frac{\rho_r(\vecv)}{\prod_{i=1}^n \rho_r(\|\widetilde{\vecb_i}\| \Z )}.
\]
Since the denominator is independent of $\vecv$, we obtain that in any given iteration,
conditioned on outputting something, the output is distributed according to
the desired distribution $D_{\Lambda+\vecc,r}$, and therefore this is also the overall
output distribution of our sampler. Moreover, by~\eqref{eq:rejsamplingprob}, 
the probability of outputting something in any given iteration is at least
$e^{-2}$, and therefore, the probability that the number of iterations is
greater than $t$ is at most $(1-e^{-2})^t$, and in particular, the expected
number of iterations is at most $e^{2}$.

It remains to show how to efficiently compute the probability in~\eqref{eq:rejsamplingprob}. By scaling, it suffices to show 
how to compute 
\[
\rho_r(\Z+c) = \sum_{k \in \Z} \exp(-\pi (k+c)^2 / r^2)
\] 
for any $r>0$ and $c \in [0,1)$. If $r<1$, the sum decays very fast, and we can achieve any desired $t$ bits of 
accuracy in time $\poly(t)$, which agrees with our notion of efficiently computing a real number (following, 
e.g., the treatment in~\cite[Section 1.4]{LovaszBook}). For $r \ge 1$, we use the Poisson summation formula (see, e.g.,~\cite[Lemma 2.8]{DBLP:journals/siamcomp/MicciancioR07}) to write
\[
\rho_r(\Z+c) = r \cdot \sum_{k \in \Z} \exp(-\pi k^2 r^2 + 2 \pi i c k) = r \cdot \sum_{k \in \Z} \exp(-\pi k^2 r^2) \cos(2 \pi c k),
\] 
which again decays fast enough so we can compute it to within any desired $t$ bits of accuracy in time $\poly(t)$.

%%% Local Variables: 
%%% mode: latex
%%% TeX-master: "lwehard"
%%% End: 

